\subsection*{把变化率逐段累积起来}
若知道一辆车每个时刻的速度，想求一段时间的净位移，可以把时间切成小段，每段用速度乘时长估计位移，再把贡献相加。速度变化时，有限切分存在误差；积分是切分变细后稳定的有向累积。

为什么叫有向？向前速度贡献正位移，倒退速度贡献负位移。净位移可能为零，但总路程并非零。同理，曲线在横轴下的部分在定积分中带负号，几何面积则需要另外处理绝对值。

\subsection*{原函数为什么能省掉求和}
若速度本来就是某个位置函数 $F$ 的导数，每小段的真实位移是两个端点的位置差。把相邻小段相加，中间端点一正一负地抵消，只剩 $F(b)-F(a)$。基本定理将这种望远镜式相消，与用导数近似小段变化的极限论证连接起来。

因此找到原函数是一种高效计算累积的方法，而不是积分唯一的定义。无法写出简单原函数时，定积分仍可能存在，也可通过数值求和逼近。反过来，形式上写出一个导数表达式，也应检查区间条件，不能跨过奇点随意用端点差。

在矩形实验中逐步加密切分，先看左端点和是否偏低，再解释原因。本例函数递增，每个矩形高度取在小段最左端，因而低于曲线。这样你不仅看到数值趋近，还知道误差的方向；这是数值计算与数学解释配合的方式。
